
\documentclass[11pt]{article}

% pdfLaTeX-friendly (no fontspec)
\usepackage[margin=1in]{geometry}
\usepackage[T1]{fontenc}
\usepackage[utf8]{inputenc}
\usepackage{lmodern}

\usepackage{amsmath,amssymb}
\usepackage{xcolor}
\usepackage{hyperref}
\usepackage{listings}

% Dirac notation (keep \rangle correct)
\newcommand{\ket}[1]{\lvert #1 \rangle}
\newcommand{\bra}[1]{\langle #1 \rvert}

% Lean listings setup (handles Lean Unicode in code blocks)
\lstdefinelanguage{Lean}{
  morekeywords={import,open,scoped,namespace,section,end,variable,def,abbrev,lemma,theorem,by,fun,have,calc,refine,intro,exact,classical,simp,funext,by_cases,omit},
  sensitive=true,
  morecomment=[l]{--},
  morecomment=[s]{/-}{-/},
  morestring=[b]"
}

\lstset{
  language=Lean,
  basicstyle=\ttfamily\footnotesize,
  columns=fullflexible,
  breaklines=true,
  frame=single,
  rulecolor=\color{black!25},
  keywordstyle=\color{blue!70!black},
  commentstyle=\color{green!40!black},
  stringstyle=\color{orange!60!black},
  showstringspaces=false,
  keepspaces=true,
  literate=
    {ℕ}{{\ensuremath{\mathbb{N}}}}1
    {𝕜}{{\ensuremath{\mathbb{k}}}}1
    {→}{{\ensuremath{\to}}}1
    {↔}{{\ensuremath{\leftrightarrow}}}1
    {∑}{{\ensuremath{\sum}}}1
    {⊆}{{\ensuremath{\subseteq}}}1
    {∈}{{\ensuremath{\in}}}1
    {∀}{{\ensuremath{\forall}}}1
    {≠}{{\ensuremath{\neq}}}1
    {↑}{{\ensuremath{\uparrow}}}1
    {·}{{\ensuremath{\bullet}}}1
    {⦃}{{\{} }1
    {⦄}{{\}} }1
    {“}{{``}}1
    {”}{{''}}1
}

\title{Explanation of the Lean 4 (mathlib) Code for the Subset-Sum (SS) Condition}
\author{}
\date{}

\begin{document}
\maketitle

\section*{Big picture}

You’re formalizing the “subset-sum support constraint” for transversal diagonal gates:

\begin{itemize}
\item A transversal diagonal gate assigns each computational basis string $s\in\{0,1\}^n$ a phase depending on a modular weighted sum
\[
\sum_j a_j s_j \bmod m.
\]
\item For a logical basis state
\[
\ket{k_L}=\sum_s c^{(k)}_s \ket{s}
\]
to be an eigenvector with eigenvalue $e^{2\pi i b_k/m}$, every string $s$ that appears with nonzero coefficient must satisfy
\[
\sum_j a_j s_j \equiv b_k \pmod m.
\]
\item You formalize that as: \textbf{support $\subseteq S_k$} (for finitary “support sets”) or as: \textbf{if amplitude is nonzero then weight $=b_k$} (for functional amplitude vectors).
\end{itemize}

Then you prove a correctness lemma:

\begin{itemize}
\item If the SS condition holds, then the abstract diagonal action \texttt{Dact} acts as a \textbf{global phase} on the state.
\end{itemize}

This is the core “SS $\Rightarrow$ eigenstate under diagonal action” statement.

\section*{Imports and namespace}

\begin{lstlisting}
import Mathlib.Data.ZMod.Basic
import Mathlib.Data.Bool.Basic
import Mathlib.Data.Finset.Lattice.Basic
import Mathlib.Algebra.BigOperators.Fin

open scoped BigOperators
namespace SS
\end{lstlisting}

\begin{itemize}
\item \texttt{ZMod.Basic}: gives \texttt{ZMod m} and basic algebra on it.
\item \texttt{Bool.Basic}: for \texttt{Bool} and simp lemmas about \texttt{if}.
\item \texttt{Finset.Lattice.Basic}: gives coercions \texttt{(↑supp : Set α)}, subset relations, etc.
\item \texttt{BigOperators.Fin}: provides \texttt{∑ i : Fin n, ...} notation.
\end{itemize}

All of your definitions are placed in \texttt{namespace SS} so names don’t pollute globally.

\section*{BitString}

\begin{lstlisting}
abbrev BitString (n : ℕ) : Type := Fin n → Bool
\end{lstlisting}

This is your encoding of a binary string of length $n$.

\begin{itemize}
\item A string $s\in\{0,1\}^n$ is a function \texttt{Fin n → Bool}.
\item This is convenient in Lean because \texttt{Fin n} is the index type $\{0,\dots,n-1\}$.
\end{itemize}

\section*{Section Basic}

\subsection*{Variables}

\begin{lstlisting}
section Basic
variable {m n : ℕ} [NeZero m]
\end{lstlisting}

Here you initially assume $m\neq 0$ so you can treat \texttt{ZMod m} as the usual modular ring.
You later discovered some lemmas don’t need it, so you used \texttt{omit} for those.

\subsection*{\texttt{bit : Bool → ZMod m}}

\begin{lstlisting}
def bit (b : Bool) : ZMod m := if b then (1 : ZMod m) else 0
\end{lstlisting}

This is the map $\{0,1\}\to\mathbb{Z}_m$ that sends
\begin{itemize}
\item \texttt{false} $\mapsto 0$
\item \texttt{true} $\mapsto 1$
\end{itemize}
This is the bridge between your Boolean bitstrings and arithmetic in \texttt{ZMod m}.

\subsection*{simp lemmas \texttt{bit\_false}, \texttt{bit\_true}}

\begin{lstlisting}
omit [NeZero m] in
@[simp] lemma bit_false : bit (m := m) false = 0 := by
  simp [bit]

omit [NeZero m] in
@[simp] lemma bit_true : bit (m := m) true = 1 := by
  simp [bit]
\end{lstlisting}

They just tell \texttt{simp} how to reduce \texttt{bit false} and \texttt{bit true}.
You wrapped them in \texttt{omit [NeZero m]} because they don’t actually use $m\neq 0$.

\subsection*{\texttt{weight}}

\begin{lstlisting}
def weight (a : Fin n → ZMod m) (s : BitString n) : ZMod m :=
  ∑ i : Fin n, a i * bit (m := m) (s i)
\end{lstlisting}

This is exactly the modular subset-sum:
\begin{itemize}
\item You provide an “angle vector” $a=(a_0,\dots,a_{n-1})$ as \texttt{Fin n → ZMod m}.
\item For a bitstring $s$, compute
\[
w_a(s)=\sum_i a_i\cdot s_i \pmod m,
\]
where $s_i\in\{0,1\}$ is implemented as \texttt{bit (s i)}.
\end{itemize}

This \texttt{weight} is the mathematical core of SS: it tells you what phase a diagonal gate assigns to $\ket{s}$.

\subsection*{\texttt{Sk} (support subset)}

\begin{lstlisting}
def Sk (a : Fin n → ZMod m) (bk : ZMod m) : Set (BitString n) :=
  { s | weight (m := m) a s = bk }
\end{lstlisting}

This is the SS subset
\[
S_k=\{\, s\in\{0,1\}^n : w_a(s)=b_k \,\}.
\]
In other words, \texttt{Sk a bk} is the set of computational basis strings that pick up the same phase label \texttt{bk}.

\subsection*{\texttt{mem\_Sk}}

\begin{lstlisting}
omit [NeZero m] in
@[simp] lemma mem_Sk {m n : ℕ} (a : Fin n → ZMod m) (bk : ZMod m) (s : BitString n) :
    s ∈ Sk (m := m) a bk ↔ weight (m := m) a s = bk :=
  Iff.rfl
\end{lstlisting}

Just the definitional unfolding of membership in that set. Again it doesn’t need \texttt{NeZero}.

\subsection*{\texttt{SSConditionSupport}}

\begin{lstlisting}
def SSConditionSupport (a : Fin n → ZMod m) (bk : ZMod m) (supp : Finset (BitString n)) : Prop :=
  (↑supp : Set (BitString n)) ⊆ Sk (m := m) a bk
\end{lstlisting}

This is the SS condition stated as a support containment statement:

\begin{quote}
Every string in the (finite) support \texttt{supp} lies inside \texttt{Sk}.
\end{quote}

This matches the paper’s statement $c^{(k)}_s\neq 0 \Rightarrow s\in S_k$,
but phrased for a \texttt{Finset} support rather than amplitudes.

\subsection*{\texttt{SSConditionSupport'}}

\begin{lstlisting}
def SSConditionSupport' (a : Fin n → ZMod m) (bk : ZMod m) (supp : Finset (BitString n)) : Prop :=
  ∀ ⦃s⦄, s ∈ supp → weight (m := m) a s = bk
\end{lstlisting}

Same condition, but in “elementwise” form:

\begin{quote}
If $s\in\texttt{supp}$, then its weight equals \texttt{bk}.
\end{quote}

\subsection*{\texttt{SSConditionSupport\_iff}}

\begin{lstlisting}
omit [NeZero m] in
theorem SSConditionSupport_iff
    {m n : ℕ} (a : Fin n → ZMod m) (bk : ZMod m) (supp : Finset (BitString n)) :
    SSConditionSupport (m := m) a bk supp ↔ SSConditionSupport' (m := m) a bk supp := by
  constructor
  · intro h s hs
    have : s ∈ (↑supp : Set (BitString n)) := by simpa using hs
    have : s ∈ Sk (m := m) a bk := h this
    simpa [Sk] using this
  · intro h s hs
    have hs' : s ∈ supp := by simpa using hs
    simpa [Sk] using h hs'

end Basic
\end{lstlisting}

This theorem proves those two formulations are logically equivalent:
\[
\texttt{SSConditionSupport a bk supp} \;\leftrightarrow\; \texttt{SSConditionSupport' a bk supp}.
\]
So you can use whichever is more convenient when proving things.

\begin{itemize}
\item The $\Rightarrow$ direction: from set inclusion you get the weight equation using \texttt{Sk}’s definition.
\item The $\Leftarrow$ direction: from the pointwise property you build the set inclusion.
\end{itemize}

Again, you \texttt{omit [NeZero m]} because this proof doesn’t need $m\neq 0$.

\section*{Section Disjointness}

\subsection*{Setup}

\begin{lstlisting}
section Disjointness

variable {m n K : ℕ}      -- removed [NeZero m]
variable (a : Fin n → ZMod m)
variable (b : Fin K → ZMod m)
variable (supp : Fin K → Finset (BitString n))
\end{lstlisting}

Here you introduce:
\begin{itemize}
\item \texttt{b : Fin K → ZMod m} = the list of intended logical phase labels $(b_0,\dots,b_{K-1})$.
\item \texttt{supp k} = the finite support set for the $k$-th logical basis state.
\end{itemize}

\subsection*{Theorem \texttt{supports\_disjoint\_of\_injective}}

\begin{lstlisting}
theorem supports_disjoint_of_injective
    (hb : Function.Injective b)
    (hSS : ∀ k, SSConditionSupport' (m := m) a (b k) (supp k)) :
    ∀ {k l : Fin K}, k ≠ l → Disjoint (supp k) (supp l) := by
  intro k l hkl
  classical
  refine Finset.disjoint_left.2 ?_
  intro s hsk hsl
  have hk : weight (m := m) a s = b k := hSS k hsk
  have hl : weight (m := m) a s = b l := hSS l hsl
  have hbl : b k = b l := by
    calc
      b k = weight (m := m) a s := hk.symm
      _   = b l               := hl
  exact hkl (hb hbl)

end Disjointness
\end{lstlisting}

Meaning:
\begin{itemize}
\item If the phase labels \texttt{b k} are all distinct (injective), and
\item each support \texttt{supp k} satisfies SS with its label \texttt{b k},
\item then supports for different logical indices cannot overlap.
\end{itemize}

This is the orthogonality-by-disjoint-support fact: if $b_k\neq b_\ell$, no basis string can appear in both
$\ket{k_L}$ and $\ket{\ell_L}$ under the SS rule.

How the proof works (in plain math):
\begin{itemize}
\item Suppose a string $s$ lies in both \texttt{supp k} and \texttt{supp l}.
\item Then SS says \texttt{weight a s = b k} and also \texttt{weight a s = b l}.
\item Hence \texttt{b k = b l}, contradicting injectivity unless $k=l$.
\item Therefore they are disjoint.
\end{itemize}

This matches the statement in your manuscript that distinct $b_k$ make $S_k$ disjoint and thus (at least) support-level orthogonality.

\section*{Section DiagonalAction}

This section switches from “supports as finsets” to “full amplitude vectors”.

\subsection*{Variables}

\begin{lstlisting}
section DiagonalAction

variable {m n : ℕ} -- no [NeZero m]
variable {𝕜 : Type*} [MonoidWithZero 𝕜]
\end{lstlisting}

\begin{itemize}
\item \texttt{𝕜} is the coefficient type for amplitudes. You only assume \texttt{MonoidWithZero} because you need multiplication (\texttt{*}) and a zero element.
\item You deliberately don’t assume \texttt{𝕜 = ℂ} yet: this keeps it abstract and reusable.
\end{itemize}

\subsection*{\texttt{State}}

\begin{lstlisting}
abbrev State (n : ℕ) (k : Type*) : Type _ := BitString n → k
\end{lstlisting}

A “state” is an amplitude assignment $\psi:\{0,1\}^n\to\mathbb{k}$.
So in physics notation: $\psi(s)$ is the coefficient of $\ket{s}$.
(You fixed the universe issue by writing \texttt{Type \_}.)

\subsection*{\texttt{SSConditionState}}

\begin{lstlisting}
def SSConditionState (a : Fin n → ZMod m) (bk : ZMod m) (ψ : State n 𝕜) : Prop :=
  ∀ s, ψ s ≠ 0 → weight (m := m) a s = bk
\end{lstlisting}

This is the most faithful Lean version of the manuscript SS condition:
if amplitude at $s$ is nonzero, then $s$ must have the right subset-sum.
It is literally the “support lies in $S_k$” condition, but formulated without explicitly defining support.

\subsection*{\texttt{Dact}}

\begin{lstlisting}
def Dact (phase : ZMod m → 𝕜) (a : Fin n → ZMod m) (ψ : State n 𝕜) : State n 𝕜 :=
  fun s => phase (weight (m := m) a s) * ψ s
\end{lstlisting}

This is an abstract diagonal operator acting on amplitude functions.

\begin{itemize}
\item In the computational basis, a diagonal gate acts by multiplying each basis vector $\ket{s}$ by some scalar depending on $s$.
\item Here, the scalar is \texttt{phase (weight a s)}.
\end{itemize}

So this corresponds exactly to the transversal diagonal gate’s induced action
\[
(D\psi)(s)=\omega^{w_a(s)}\,\psi(s),
\]
but abstracted as \texttt{phase : ZMod m → 𝕜} instead of hardcoding $\omega^x$.

\subsection*{Theorem \texttt{Dact\_eq\_global\_of\_SS}}

\begin{lstlisting}
theorem Dact_eq_global_of_SS
    (phase : ZMod m → 𝕜) (a : Fin n → ZMod m) (bk : ZMod m) (ψ : State n 𝕜)
    (hSS : SSConditionState (m := m) a bk ψ) :
    Dact (m := m) phase a ψ = fun s => phase bk * ψ s := by
  classical
  funext s
  by_cases hs : ψ s = 0
  · simp [Dact, hs]
  · have hw : weight (m := m) a s = bk := hSS s hs
    simp [Dact, hw]

end DiagonalAction

end SS
\end{lstlisting}

This is the “SS correctness” lemma:
If $\psi$ only has support on strings with weight $=b_k$, then applying the diagonal action multiplies every nonzero component by the same scalar \texttt{phase bk}.
So the whole state is an eigenvector with eigenvalue \texttt{phase bk}.

How the proof works:
\begin{itemize}
\item Prove equality of functions, so \texttt{funext s}.
\item Split into two cases:
\begin{enumerate}
\item $\psi(s)=0$: both sides are $0$ after multiplication, so \texttt{simp} closes it.
\item $\psi(s)\neq 0$: SS tells you \texttt{weight a s = bk}, so the phase is \texttt{phase bk}, and \texttt{simp} closes it.
\end{enumerate}
\end{itemize}

This is the formal Lean statement that SS implies the state is an eigenstate of the diagonal gate.

\section*{What this file verifies (and what it doesn’t yet)}

\subsection*{Verified / formalized}

\begin{itemize}
\item The definition of subset-sum weights on bitstrings.
\item The definition of the SS subsets $S_k$.
\item The SS condition as “support subset” and as “nonzero amplitude $\Rightarrow$ correct congruence”.
\item SS $\Rightarrow$ disjoint supports (under distinct phase labels).
\item SS $\Rightarrow$ eigenstate under the abstract diagonal action (global phase).
\end{itemize}

\end{document}
